% ASSERT_CONVENTION: natural_units=natural, metric_signature=mostly_minus, fourier_convention=physics, coupling_convention=alpha_s, generator_normalization=T(fund)=1/2, covariant_derivative_sign=D_mu=partial_mu-igA
%
% Lecture 7: The Dualised Standard Model
% Part of: The Dualised Standard Model -- Part III Lecture Notes
%
% Conventions (Seiberg hep-th/9411149, unchanged from Lectures 1-6):
%   T(fund) = 1/2 per Weyl fermion
%   b_0 = (11N_c - 2N_f)/3 for N_f Dirac flavours
%   A(fund) = 1 (cubic anomaly coefficient per fundamental Weyl fermion)
%   alpha_s = g_s^2 / (4 pi)
%   R[Q] = (N_f - N_c)/N_f; R[W] = 2
%   N_f counts Dirac flavour pairs (each = 2 Weyl fermions)
%
% NEW CONVENTIONS FOR THIS LECTURE:
%   Hypercharge: Y = T_R^3 + Q_V/6  (Cacciapaglia-Sannino 2407.17281 Eq. 6)
%   Epistemic: all non-SUSY duality claims carry conjectural qualifiers (P1)
%
\documentclass[11pt,a4paper]{article}
% ASSERT_CONVENTION: natural_units=natural, metric_signature=mostly_minus, fourier_convention=physics, coupling_convention=alpha_s, generator_normalization=T(fund)=1/2, covariant_derivative_sign=D_mu=partial_mu-igA
%
% CONVENTIONS (Seiberg hep-th/9411149)
% Metric: (+,-,-,-) (mostly minus)
% T(fund) = 1/2 per Weyl fermion
% b_0 = (11N_c - 2N_f)/3 for N_f Dirac flavours
% D_mu = partial_mu - ig A_mu^a T^a
% A(fund) = 1 (cubic anomaly coefficient per fundamental Weyl fermion)
% alpha_s = g_s^2 / (4 pi)
% R[Q] = (N_f - N_c)/N_f; R[W] = 2
% N_f counts Dirac flavour pairs (each = 2 Weyl fermions)
%
% This file is \input{} by all lecture files.

%% ---------- Document class ----------
% (loaded by the wrapper; preamble only provides packages and macros)

%% ---------- Standard packages ----------
\usepackage{amsmath,amssymb,amsthm,mathtools}
\usepackage{hyperref}
\usepackage[capitalise,noabbrev]{cleveref}
\usepackage{enumitem}
\usepackage{booktabs}
\usepackage{array}
\usepackage{xcolor}
\usepackage{tikz}
\usetikzlibrary{arrows.meta,decorations.markings,positioning,calc}
\usepackage{fancybox}
\usepackage{tcolorbox}
\tcbuselibrary{skins,breakable}
\usepackage{slashed}              % Feynman slash notation
\usepackage{cite}                 % compressed citation lists

%% ---------- Hyperref configuration ----------
\hypersetup{
  colorlinks=true,
  linkcolor=blue!70!black,
  citecolor=green!50!black,
  urlcolor=blue!80!black,
  pdfauthor={Part III Lecture Notes},
  pdftitle={The Dualised Standard Model}
}

%% ---------- Theorem environments ----------
\theoremstyle{plain}
\newtheorem{theorem}{Theorem}[section]
\newtheorem{lemma}[theorem]{Lemma}
\newtheorem{proposition}[theorem]{Proposition}
\newtheorem{corollary}[theorem]{Corollary}

\theoremstyle{definition}
\newtheorem{definition}[theorem]{Definition}
\newtheorem{example}[theorem]{Example}
\newtheorem{exercise}{Exercise}[section]

\theoremstyle{remark}
\newtheorem{remark}[theorem]{Remark}

%% ---------- Physics macros: gauge groups and representations ----------
\newcommand{\Nc}{N_{\!c}}
\newcommand{\Nf}{N_{\!f}}
\newcommand{\SU}[1]{\mathrm{SU}(#1)}
\newcommand{\UU}[1]{\mathrm{U}(#1)}

% Representations
\newcommand{\fund}{\square}
\newcommand{\antifund}{\overline{\square}}
\newcommand{\adj}{\mathrm{adj}}
\newcommand{\antisym}{\wedge^{\!2}}        % antisymmetric rank-2
\newcommand{\sym}{\mathrm{S}_2}           % symmetric rank-2

%% ---------- Physics macros: group theory constants ----------
% Dynkin index: Tr(T^a T^b) = T(R) delta^{ab}
\newcommand{\Tfund}{\tfrac{1}{2}}          % T(fund) = 1/2 per Weyl fermion
\newcommand{\Tadj}{\Nc}                    % T(adj) = N_c
\newcommand{\Cfund}{\frac{\Nc^2 - 1}{2\Nc}}  % C_2(fund)
\newcommand{\Cadj}{\Nc}                    % C_2(adj)

% Cubic anomaly coefficient: A(R) = Tr_R[T^a {T^b, T^c}] with A(fund) = 1
\newcommand{\Afund}{1}

%% ---------- Physics macros: beta function ----------
\newcommand{\bzero}{b_0}
\newcommand{\bone}{b_1}
\newcommand{\alphas}{\alpha_s}
\newcommand{\gs}{g_s}

%% ---------- Physics macros: SUSY / R-charge ----------
\newcommand{\Rcharge}[1]{R[#1]}
\newcommand{\Wsup}{W}                      % superpotential

%% ---------- Physics macros: covariant derivative and field strength ----------
\newcommand{\Dmu}{D_\mu}
\newcommand{\Fmn}{F_{\mu\nu}}
\newcommand{\Ftilde}{\widetilde{F}_{\mu\nu}}

%% ---------- Physics macros: common abbreviations ----------
\newcommand{\ie}{\textit{i.e.}}
\newcommand{\eg}{\textit{e.g.}}
\newcommand{\cf}{\textit{cf.}}
\newcommand{\IR}{\ensuremath{\mathrm{IR}}}
\newcommand{\UV}{\ensuremath{\mathrm{UV}}}
\newcommand{\AF}{\ensuremath{\mathrm{AF}}}
\newcommand{\BZ}{\ensuremath{\mathrm{BZ}}}
\newcommand{\NSVZ}{\ensuremath{\mathrm{NSVZ}}}
\newcommand{\SQCD}{\ensuremath{\mathrm{SQCD}}}
\newcommand{\QCD}{\ensuremath{\mathrm{QCD}}}
\newcommand{\SM}{\ensuremath{\mathrm{SM}}}
\newcommand{\Tr}{\mathrm{Tr}}

%% ---------- Convention box macro ----------
\newtcolorbox{conventionbox}[1][]{%
  enhanced,
  colback=blue!5!white,
  colframe=blue!60!black,
  fonttitle=\bfseries,
  title={Conventions},
  #1
}

%% ---------- Comparison table style ----------
\newtcolorbox{comparisontable}[1][]{%
  enhanced,
  colback=orange!5!white,
  colframe=orange!70!black,
  fonttitle=\bfseries,
  title={Comparison},
  #1
}

%% ---------- Warning / important box ----------
\newtcolorbox{warningbox}[1][]{%
  enhanced,
  colback=red!5!white,
  colframe=red!70!black,
  fonttitle=\bfseries,
  title={Important},
  #1
}

%% ---------- Preview / forward-reference box ----------
\newtcolorbox{previewbox}[1][]{%
  enhanced,
  colback=green!5!white,
  colframe=green!50!black,
  fonttitle=\bfseries,
  title={Preview},
  #1
}

%% ---------- Page layout ----------
\usepackage[margin=2.5cm]{geometry}
\setlength{\parskip}{0.4em}
\setlength{\parindent}{0em}

%% ---------- Section formatting ----------
\usepackage{titlesec}
\titleformat{\section}{\Large\bfseries}{\thesection}{1em}{}
\titleformat{\subsection}{\large\bfseries}{\thesubsection}{1em}{}

%% ---------- End of preamble ----------


\title{\textbf{Lecture 7: The Dualised Standard Model}}
\author{The Dualised Standard Model --- Part III Lecture Notes}
\date{Lent Term 2027}

\begin{document}
\maketitle

\tableofcontents
\newpage

%% ========================================================================
%%  CONVENTIONS
%% ========================================================================

\begin{conventionbox}[title={Conventions for Lecture 7}]
All conventions from Lectures~1--6 remain in force.  New conventions
for the Pati--Salam embedding and hypercharge identification are
introduced below.

\renewcommand{\arraystretch}{1.3}
\begin{tabular}{@{}l l l@{}}
\toprule
\textbf{Quantity} & \textbf{Convention} & \textbf{Comment} \\
\midrule
Units & $\hbar = c = 1$ (natural) & \\
Metric & $\eta_{\mu\nu} = \mathrm{diag}(+1,-1,-1,-1)$ & mostly minus \\
Dynkin index & $\Tr(T^a T^b) = \Tfund\, \delta^{ab}$ & fundamental of $\SU{\Nc}$ \\
Cubic anomaly & $A(\fund) = 1$, $A(\antifund) = -1$ & per fundamental Weyl \\
Beta function & $\bzero = (11\Nc - 2\Nf)/3$ & $\Nf$ Dirac flavours \\
$\Nf$ counting & Dirac flavour pairs & each $= 2$ Weyl fermions \\
\midrule
\multicolumn{3}{@{}l}{\textbf{New in this lecture:}} \\
Hypercharge & $Y = T_R^3 + \tfrac{1}{6}\,Q_V$ & Eq.~(6) of \cite{CacciapagliaSannino2024} \\
Duality map & $X = \Nf - N$ & magnetic gauge group $\SU{X}$ \\
Epistemic & ``if the duality holds'' & all non-SUSY claims (P1) \\
\bottomrule
\end{tabular}

\medskip

\textbf{Epistemic convention.}  Every statement asserting or assuming
non-supersymmetric duality carries an explicit conjectural qualifier.
The field content \emph{tables} and anomaly matching \emph{equations}
are exact group-theoretic results; the \emph{duality interpretation}
is conjectural.
\end{conventionbox}


\begin{center}
\textit{Gauge group roadmap:} \\[4pt]
\S\ref{sec:general-duality}: General $\SU{N}$ duality ansatz $\;\to\;$
\S\ref{sec:pati-salam}: Pati--Salam $\SU{4}_{PS}$ embedding $\;\to\;$
\S\ref{sec:magnetic-spectrum}: Magnetic spectrum for $\SU{3}_c \times \SU{2}_L \times \UU{1}_Y$ $\;\to\;$\\[2pt]
\S\ref{sec:electric-theory}: Electric $\SU{2}$ theory $\;\to\;$
\S\ref{sec:anomaly-matching}: Anomaly matching $\;\to\;$
\S\ref{sec:summary-lec7}: Summary
\end{center}

%% ========================================================================
%%  SECTION 1: The General Duality Framework
%% ========================================================================
\section{The General Duality Framework}
\label{sec:general-duality}

In Lecture~6, we introduced the conjectured non-SUSY duality programme
and stated its central vulnerability: anomaly matching is the sole
surviving consistency check.  We now \emph{apply} this framework to
construct the dualised Standard Model, following Cacciapaglia and
Sannino~\cite{CacciapagliaSannino2024}.

\subsection{The electric theory}

The starting point is an $\SU{N}$ gauge theory with the following
matter content~\cite{CacciapagliaSannino2024}:
\begin{itemize}
  \item $\Nf$ Dirac fermions $Q^i$, $\widetilde{Q}_i$
    ($i = 1,\ldots,\Nf$) in the fundamental and antifundamental
    representations of $\SU{N}$;
  \item One Weyl fermion $\lambda$ in the adjoint representation of
    $\SU{N}$.
\end{itemize}
This is the \emph{non-SUSY cousin} of $\mathcal{N}=1$ SQCD with an
adjoint: the same fermion content as the SUSY theory, but without
squarks or sleptons as fundamental fields.  The quantum global symmetry
group is:
\begin{equation}\label{eq:electric-global}
  G_{\mathrm{global}}
  \;=\; \SU{\Nf}_L \times \SU{\Nf}_R \times \UU{1}_V
  \times \UU{1}_{AF}\,,
\end{equation}
where $\UU{1}_V$ is the vector (baryon-number) symmetry acting as
$Q \to e^{i\alpha} Q$, $\widetilde{Q} \to e^{-i\alpha}\widetilde{Q}$,
and $\UU{1}_{AF}$ is an anomaly-free axial symmetry under which the
adjoint fermion $\lambda$ is charged.

\begin{center}
\renewcommand{\arraystretch}{1.3}
\begin{tabular}{c c c c c c}
\toprule
\multicolumn{6}{c}{\textbf{Electric theory (UV)}} \\
\midrule
Fields & $\SU{N}$ & $\SU{\Nf}_L$ & $\SU{\Nf}_R$ & $\UU{1}_V$ & $\UU{1}_{AF}$ \\
\midrule
$\lambda$ & Adj & $\mathbf{1}$ & $\mathbf{1}$ & $0$ & $1$ \\
$Q$       & $F$ & $F$ & $\mathbf{1}$ & $1$  & $-N/\Nf$ \\
$\widetilde{Q}$ & $\bar{F}$ & $\mathbf{1}$ & $\bar{F}$ & $-1$ & $-N/\Nf$ \\
\bottomrule
\end{tabular}
\captionof{table}{Electric theory field content
(Table~I, top block, of~\cite{CacciapagliaSannino2024}).
$F$ denotes the fundamental representation, $\bar{F}$ the
antifundamental, and Adj the adjoint.}
\label{tab:electric-general}
\end{center}

\medskip

\noindent\textbf{Key feature:} The electric theory contains
\textbf{no elementary scalars}.  All gauge fields and matter fields are
fermions (and gauge bosons).  \emph{If the duality holds}, all scalars
in the low-energy description are composites.

\subsection{The conjectured magnetic dual}

\emph{If the conjectured non-SUSY duality of Lecture~6 holds}, the
magnetic (IR) description is an $\SU{X}$ gauge theory with~\cite{CacciapagliaSannino2024}
\begin{equation}\label{eq:duality-map}
  \boxed{X \;=\; \Nf - N\,,}
\end{equation}
and the following quasi-supersymmetric spectrum (Table~I, bottom block,
of~\cite{CacciapagliaSannino2024}).  The term ``quasi-supersymmetric''
reflects the organisation of the magnetic degrees of freedom into
multiplet-like pairs reminiscent of SUSY, even though \emph{no}
supersymmetry is imposed~\cite{AMPS2011}:
\begin{equation}\label{eq:quasi-susy}
  (G_\mu,\, \lambda_m)\,,
  \qquad
  (q,\, \phi)\,,
  \qquad
  (\widetilde{q},\, \widetilde{\phi})\,,
  \qquad
  (M,\, \Phi_H)\,.
\end{equation}
Each pair consists of a fermion and a scalar partner.  The complete
quantum numbers are:

\begin{center}
\renewcommand{\arraystretch}{1.3}
\begin{tabular}{c c c c c c}
\toprule
\multicolumn{6}{c}{\textbf{Magnetic theory (IR)}} \\
\midrule
Fields & $\SU{X}$ & $\SU{\Nf}_L$ & $\SU{\Nf}_R$ & $\UU{1}_V$ & $\UU{1}_{AF}$ \\
\midrule
$\lambda_m$ (gaugino) & Adj & $\mathbf{1}$ & $\mathbf{1}$ & $0$ & $1$ \\
$q$ (mag.\ quark) & $F$ & $\bar{F}$ & $\mathbf{1}$ & $N/X$  & $-X/\Nf$ \\
$\widetilde{q}$ (mag.\ antiquark) & $\bar{F}$ & $\mathbf{1}$ & $F$ & $-N/X$ & $-X/\Nf$ \\
$M$ (mesino) & $\mathbf{1}$ & $F$ & $\bar{F}$ & $0$ & $-1+2X/\Nf$ \\
$\phi$ (squark) & $F$ & $\bar{F}$ & $\mathbf{1}$ & $N/X$ & $1-X/\Nf$ \\
$\widetilde{\phi}$ (anti-squark) & $\bar{F}$ & $\mathbf{1}$ & $F$ & $-N/X$ & $1-X/\Nf$ \\
$\Phi_H$ (mes-Higgs) & $\mathbf{1}$ & $F$ & $\bar{F}$ & $0$ & $2X/\Nf$ \\
\bottomrule
\end{tabular}
\captionof{table}{Magnetic theory field content
(Table~I, bottom block, of~\cite{CacciapagliaSannino2024}).
\textbf{Fermions:} $\lambda_m$, $q$, $\widetilde{q}$, $M$.
\textbf{Scalars:} $\phi$, $\widetilde{\phi}$, $\Phi_H$.
The naming reflects the quasi-SUSY pairing~\eqref{eq:quasi-susy}.}
\label{tab:magnetic-general}
\end{center}

\medskip

The \emph{magnetic interactions}, \emph{if the duality holds}, take the
form~\cite{CacciapagliaSannino2024}:
\begin{equation}\label{eq:mag-lagrangian}
  \mathcal{L}_m \;\supset\;
  y\, q\,\widetilde{q}\, \Phi_H
  \;+\; y'\, q\,\widetilde{\phi}\, M
  \;+\; \widetilde{y}'\, \widetilde{q}\, \phi\, M
  \;+\; \xi_L\, \lambda_m\, q\, \phi^\dagger
  \;+\; \xi_R\, \lambda_m\, \widetilde{q}\, \widetilde{\phi}^\dagger
  \;+\; \text{h.c.}
  \;+\; \text{scalar quartics}\,,
\end{equation}
where the first three terms emerge from a superpotential coupling if
$y = y' = \widetilde{y}'$, and the last two from gauge interactions if
$\xi_L = \xi_R = g_m$ (Eq.~(4) of~\cite{CacciapagliaSannino2024}).
The \textbf{universal Yukawa coupling} $y$ is a central prediction of
the framework.

\subsection{Composite fields as electric bilinears}

The scalar fields in the magnetic theory are naturally interpreted as
composite states of the electric theory~\cite{CacciapagliaSannino2024}:
\begin{equation}\label{eq:composites}
  M \;=\; \langle \widetilde{Q}\, \lambda\, Q \rangle\,,
  \qquad
  \Phi_H \;=\; \langle \widetilde{Q}\, \lambda\, \lambda\, Q \rangle\,.
\end{equation}
The mesino $M$ is a three-fermion composite, while the mes-Higgs
$\Phi_H$ is a four-fermion composite.  \textbf{If the duality holds, the Higgs field is not
an elementary scalar}---within the conjectured dual map, it is a deeply composite state of the electric
theory.  This is the non-SUSY analogue of the Seiberg duality meson
$M \sim Q\widetilde{Q}$ from Lecture~3, but with the adjoint fermion
$\lambda$ providing additional ``glue.''

\begin{remark}[Connection to Seiberg duality]
  Compare with Lecture~3: in SUSY Seiberg duality, the meson
  $M^i_j = Q^i \widetilde{Q}_j / \Lambda$ is a two-fermion composite.
  The non-SUSY version requires the adjoint fermion $\lambda$ to
  ``saturate'' the additional quantum numbers.  The quasi-SUSY
  multiplet structure~\eqref{eq:quasi-susy} mirrors the SUSY
  structure (vector multiplet, chiral multiplets, meson multiplet)
  even though no supersymmetry is imposed.
\end{remark}


%% ========================================================================
%%  SECTION 2: Pati-Salam Embedding and Hypercharge
%% ========================================================================
\section{Pati--Salam Embedding and Hypercharge}
\label{sec:pati-salam}

The general duality framework of Section~\ref{sec:general-duality}
applies to any $\SU{N}$ gauge theory.  To connect it to the Standard
Model, we must embed the SM gauge group
$\SU{3}_c \times \SU{2}_L \times \UU{1}_Y$ within the duality
structure.  Following Sannino~\cite{Sannino2011} and Cacciapaglia and
Sannino~\cite{CacciapagliaSannino2024}, we use the Pati--Salam
embedding~\cite{PS1974}.

\subsection{Lepton number as the fourth colour}

Pati and Salam~\cite{PS1974} observed that quarks and leptons can be
unified by extending the colour group $\SU{3}_c$ to $\SU{4}$, with the
lepton as the ``fourth colour'':
\begin{equation}\label{eq:pati-salam}
  \psi_C \;=\; (u, d, s, \ldots, \nu_e, e, \ldots)_C\,,
  \qquad C = 1,2,3,4\,,
\end{equation}
where $C = 1,2,3$ are the ordinary QCD colours and $C = 4$ is the
``lepton colour.''  Under the decomposition
$\SU{4} \to \SU{3}_c \times \UU{1}_{B-L}$, the fundamental
$\mathbf{4}$ splits as:
\begin{equation}\label{eq:su4-decomp}
  \mathbf{4} \;\to\; \mathbf{3}_{+1/3} \;\oplus\; \mathbf{1}_{-1}\,,
\end{equation}
where the subscript is the $B-L$ charge (baryon number minus lepton
number), normalised so that quarks have $B-L = +1/3$ and leptons have
$B-L = -1$.

\subsection{Flavour symmetry decomposition}

The flavour symmetry of the electric theory~\eqref{eq:electric-global}
contains $\SU{\Nf}_{L/R}$.  For the SM embedding, we decompose the
flavour group as~\cite{CacciapagliaSannino2024}:
\begin{equation}\label{eq:flavour-decomp}
  \SU{\Nf}_{L/R}
  \;\supset\;
  \SU{n_g} \times \SU{2}_{L/R}\,,
\end{equation}
where $n_g = \Nf / 2$ is the number of SM generations and $\SU{2}_R$
is (part of) the right-handed isospin symmetry.  For three generations,
$n_g = 3$ and $\Nf = 2 n_g = 6$.

The electroweak gauge group $\SU{2}_L$ is identified with the
left-handed factor, and $\SU{2}_R$ provides the quantum number
$T_R^3$ needed to define hypercharge.

\subsection{Hypercharge identification}

The SM hypercharge is defined as~\cite{CacciapagliaSannino2024}:
\begin{equation}\label{eq:hypercharge}
  \boxed{Y \;=\; T_R^3 \;+\; \frac{1}{6}\, Q_V\,,}
\end{equation}
where $T_R^3$ is the third component of $\SU{2}_R$ isospin and $Q_V$
is the $\UU{1}_V$ charge from~\eqref{eq:electric-global}.

\begin{warningbox}[title={Hypercharge check: reproducing SM values}]
  Formula~\eqref{eq:hypercharge} \emph{must} reproduce the standard SM
  hypercharges.  We verify this explicitly.

  \medskip

  Under $\SU{2}_L \times \SU{2}_R$, the SM fermions are organised as
  left-handed doublets ($T_R^3 = 0$) and right-handed singlets
  ($T_R^3 = \pm\tfrac{1}{2}$).  The $Q_V$ charge equals
  $+1$ for quarks ($Q$) and $-3$ for leptons added as elementary
  states for anomaly cancellation.

  \renewcommand{\arraystretch}{1.4}
  \begin{center}
  \begin{tabular}{l c c c c}
  \toprule
  \textbf{SM field} & $T_R^3$ & $Q_V$ & $Y = T_R^3 + Q_V/6$
    & \textbf{SM value} \\
  \midrule
  $q_L = (u_L, d_L)$ & $0$ & $+1$ & $0 + 1/6 = +1/6$ & $+1/6$ \;\checkmark \\
  $u_R$ & $+1/2$ & $+1$ & $1/2 + 1/6 = +2/3$ & $+2/3$ \;\checkmark \\
  $d_R$ & $-1/2$ & $+1$ & $-1/2 + 1/6 = -1/3$ & $-1/3$ \;\checkmark \\
  $\ell_L = (\nu_L, e_L)$ & $0$ & $-3$ & $0 - 3/6 = -1/2$ & $-1/2$ \;\checkmark \\
  $\nu_R$ & $+1/2$ & $-3$ & $1/2 - 3/6 = 0$ & $0$ \;\checkmark \\
  $e_R$ & $-1/2$ & $-3$ & $-1/2 - 3/6 = -1$ & $-1$ \;\checkmark \\
  \bottomrule
  \end{tabular}
  \end{center}

  \smallskip

  All six SM hypercharges are reproduced exactly.  The
  relation to the alternative convention $Y = T_R^3 + (B-L)/2$
  is $Q_V = 3(B-L)$.
\end{warningbox}

\subsection{The generation bound}
\label{ssec:generation-bound}

The electric gauge group is $\SU{2n_g - 4}$, where $n_g$ is the number
of generations~\cite{Sannino2011,CacciapagliaSannino2024}.  For this
group to be non-abelian (a necessary condition for a non-trivial
gauge theory), we require:
\begin{equation}\label{eq:gen-bound}
  2n_g - 4 \;\geq\; 2
  \qquad \Longrightarrow \qquad
  \boxed{n_g \;\geq\; 3\,.}
\end{equation}
\emph{If the duality holds}, the existence of at least three
generations of quarks and leptons is \textbf{not an accident}: it is a
\textbf{mathematical requirement} for the electric dual to be a
non-trivial gauge theory~\cite{Sannino2011}.

For $n_g = 3$ (three generations):
\begin{itemize}
  \item Electric gauge group: $\SU{2 \cdot 3 - 4} = \SU{2}$.
  \item Number of Dirac flavours: $\Nf = 2 n_g = 6$.
  \item Magnetic gauge group (from QCD duality): $X = \Nf - N = 6 - 3 = 3$,
    but only $\SU{3}_c$ (the QCD subgroup of $\SU{4}_{\mathrm{PS}}$)
    is gauged in the SM.  The full Pati--Salam $\SU{4}$ is the
    electric theory containing both quarks and leptons.
\end{itemize}

\begin{remark}[Upper bound on generations]
  Requiring asymptotic freedom of the magnetic theory further
  constrains $n_g \leq 6$~\cite{Sannino2011}.  The observed value
  $n_g = 3$ is the \emph{minimal} number of generations consistent
  with the duality framework.
\end{remark}


%% ========================================================================
%%  SECTION 3: The Magnetic Spectrum -- SM Fields and Quasi-SUSY Partners
%% ========================================================================
\section{The Magnetic Spectrum: SM Fields and Quasi-SUSY Partners}
\label{sec:magnetic-spectrum}

We now specialise the general magnetic field content of
Table~\ref{tab:magnetic-general} to the SM values $N = 3$, $\Nf = 6$
(three generations), giving $X = \Nf - N = 3$.
Following Table~II of~\cite{CacciapagliaSannino2024}, the full
magnetic spectrum under $\SU{3}_c \times \SU{6}_L \times \SU{6}_R
\times \UU{1}_V \times \UU{1}_{AF}$ is:

\begin{center}
\renewcommand{\arraystretch}{1.3}
\begin{tabular}{c c c c c c}
\toprule
\multicolumn{6}{c}{\textbf{Electric theory (UV): $\SU{3}$ with $\Nf = 6$}} \\
\midrule
Fields & $\SU{3}$ & $\SU{6}_L$ & $\SU{6}_R$ & $\UU{1}_V$ & $\UU{1}_{AF}$ \\
\midrule
$\lambda$ & Adj & $\mathbf{1}$ & $\mathbf{1}$ & $0$ & $1$ \\
$Q$       & $F$ & $F$ & $\mathbf{1}$ & $1$  & $-1/2$ \\
$\widetilde{Q}$ & $\bar{F}$ & $\mathbf{1}$ & $\bar{F}$ & $-1$ & $-1/2$ \\
\midrule
\multicolumn{6}{c}{\textbf{+ Elementary states for anomaly cancellation:}} \\
$L$   & $\mathbf{1}$ & $F$ & $\mathbf{1}$ & $-3$ & $0$ \\
$\widetilde{L}$ & $\mathbf{1}$ & $\mathbf{1}$ & $\bar{F}$ & $3$ & $0$ \\
\bottomrule
\end{tabular}
\captionof{table}{Electric theory field content for $N=3$, $\Nf=6$
(Table~II, top block, of~\cite{CacciapagliaSannino2024}).
The lepton fields $L$, $\widetilde{L}$ are elementary states
required for $\UU{1}_Y$ anomaly cancellation.}
\label{tab:electric-SM}
\end{center}

\begin{center}
\renewcommand{\arraystretch}{1.3}
\begin{tabular}{c c c c c c}
\toprule
\multicolumn{6}{c}{\textbf{Magnetic theory (IR): $\SU{3}_c$ with $\Nf = 6$}} \\
\midrule
Fields & $\SU{3}_c$ & $\SU{6}_L$ & $\SU{6}_R$ & $\UU{1}_V$ & $\UU{1}_{AF}$ \\
\midrule
$\lambda_m$ (gaugino) & Adj & $\mathbf{1}$ & $\mathbf{1}$ & $0$ & $1$ \\
$q$ (mag.\ quarks)    & $F$ & $\bar{F}$ & $\mathbf{1}$ & $1$  & $-1/2$ \\
$\widetilde{q}$ (mag.\ antiquarks) & $\bar{F}$ & $\mathbf{1}$ & $F$ & $-1$ & $-1/2$ \\
$l \equiv L$ (leptons)  & $\mathbf{1}$ & $F$ & $\mathbf{1}$ & $-3$ & $0$ \\
$\widetilde{l} \equiv \widetilde{L}$ (antileptons) & $\mathbf{1}$ & $\mathbf{1}$ & $\bar{F}$ & $3$ & $0$ \\
$M$ (mesino)  & $\mathbf{1}$ & $F$ & $\bar{F}$ & $0$ & $0$ \\
$\phi$ (squarks) & $F$ & $\bar{F}$ & $\mathbf{1}$ & $1$ & $1/2$ \\
$\widetilde{\phi}$ (anti-squarks) & $\bar{F}$ & $\mathbf{1}$ & $F$ & $-1$ & $1/2$ \\
$\Phi_H$ (mes-Higgs) & $\mathbf{1}$ & $F$ & $\bar{F}$ & $0$ & $1$ \\
\bottomrule
\end{tabular}
\captionof{table}{Magnetic theory field content for $N=3$, $\Nf=6$
(Table~II, bottom block, of~\cite{CacciapagliaSannino2024}).
\textbf{Fermions:} $\lambda_m$, $q$, $\widetilde{q}$, $l$, $\widetilde{l}$, $M$.
\textbf{Scalars:} $\phi$, $\widetilde{\phi}$, $\Phi_H$.}
\label{tab:magnetic-SM}
\end{center}

\subsection{Identification of SM fields}

The magnetic spectrum of Table~\ref{tab:magnetic-SM} contains
\textbf{all SM fields plus additional quasi-SUSY partners}:

\paragraph{SM quarks and leptons.}
The magnetic quarks $q$, $\widetilde{q}$ are, if the duality holds, identified with the SM
quarks.  Under the flavour decomposition~\eqref{eq:flavour-decomp},
each $q$ ($\bar{F}$ of $\SU{6}_L$) decomposes into an $\SU{3}_{\mathrm{gen}}$
triplet (three generations) of $\SU{2}_L$ doublets and singlets,
reproducing the SM quark content.  The leptons $l \equiv L$ and
$\widetilde{l} \equiv \widetilde{L}$ are elementary states---they are
\textbf{not} composites of the electric theory in the pure QCD duality.

\begin{warningbox}[title={Leptons in the QCD dual vs.\ Pati--Salam}]
  In the \textbf{QCD-sector duality} (Table~\ref{tab:magnetic-SM}),
  leptons $l$, $\widetilde{l}$ are \textbf{elementary}
  states added for $\UU{1}_Y$ anomaly cancellation.  They are colour
  singlets and do not participate in the $\SU{3}_c$ dynamics.

  In the \textbf{full Pati--Salam extension}~\cite{Sannino2011},
  leptons are the $C = 4$ component of $\SU{4}$ fundamentals, and the
  distinction between quarks and leptons is a consequence of
  $\SU{4} \to \SU{3}_c \times \UU{1}_{B-L}$ breaking.

  This lecture follows the QCD-sector duality of~\cite{CacciapagliaSannino2024}, where leptons are elementary.
\end{warningbox}

\paragraph{The 18 Higgs doublets.}
The mes-Higgs field $\Phi_H$ is a gauge singlet that transforms as
$(\mathbf{6}, \bar{\mathbf{6}})$ under $\SU{6}_L \times \SU{6}_R$.
Under the flavour decomposition~\eqref{eq:flavour-decomp}, $\Phi_H$
decomposes into a matrix of bi-doublets:
\begin{equation}\label{eq:Phi-H}
  \Phi_H \;=\; \{H_{ij}\}\,,
  \qquad i,j = 1,2,3\,,
\end{equation}
where the indices $i,j$ label the three generations
(from $\SU{3}_{\mathrm{gen}}$) and each $H_{ij}$ is a complex bi-doublet
of $\SU{2}_L \times \SU{2}_R$.  Each bi-doublet decomposes into two
$\SU{2}_L$ doublets distinguished by hypercharge~\cite{CacciapagliaSannino2024}:
\begin{equation}\label{eq:Higgs-decomp}
  H_{ij} \;=\; (H_{ij}^u,\; H_{ij}^d)\,,
\end{equation}
where $H_{ij}^u$ has $Y = +1/2$ and $H_{ij}^d$ has $Y = -1/2$.
The total Higgs content is therefore:
\begin{equation}\label{eq:18-Higgs}
  \underbrace{9 \text{ bi-doublets}}_{i,j = 1,2,3}
  \;\times\;
  \underbrace{2 \text{ doublets per bi-doublet}}_{H^u,\, H^d}
  \;=\; \boxed{18 \text{ Higgs doublets}\,.}
\end{equation}
One of these---$H_0 \equiv H_{33}^u$---acquires the
electroweak VEV $\langle H_0 \rangle \approx v$; the remaining 17 have
positive mass-squared at tree level (they are ``dormant'' Higgs
doublets).

\paragraph{Quasi-SUSY partners.}
The magnetic spectrum includes states \emph{beyond} the minimal SM
content:
\begin{itemize}
  \item \textbf{Gaugino} $\lambda_m$: a colour-octet Weyl fermion
    (Adj of $\SU{3}_c$), the quasi-SUSY partner of the gluon.
  \item \textbf{Squarks} $\phi$, $\widetilde{\phi}$: colour-triplet
    and anti-triplet \emph{scalars}, the quasi-SUSY partners of
    $q$ and $\widetilde{q}$.
  \item \textbf{Mesino} $M$: a colour-singlet fermion transforming as
    $(\mathbf{6}, \bar{\mathbf{6}})$ under $\SU{6}_L \times \SU{6}_R$,
    the quasi-SUSY partner of $\Phi_H$.
\end{itemize}
\emph{If the duality holds}, these states are expected to appear at the
scale $\Lambda_S \sim \text{few TeV}$, where the quasi-supersymmetry
of the magnetic theory becomes
manifest~\cite{CacciapagliaSannino2024}.  Their masses are
generically required to be at or above collider bounds (the gaugino
and squarks are constrained by LHC searches to be above roughly
$2\;\text{TeV}$~\cite{CacciapagliaSannino2024}).

\subsection{The universal Yukawa coupling}

The magnetic Lagrangian~\eqref{eq:mag-lagrangian} generates a
\textbf{single universal Yukawa coupling} $y$ between the magnetic
quarks and the mes-Higgs field.  In terms of the individual Higgs
doublets~\cite{CacciapagliaSannino2024}:
\begin{equation}\label{eq:yukawa}
  y\, q\, \widetilde{q}\, \Phi_H
  \;=\; y \sum_{i,j}
  \Bigl(
    q_L^i\, u_R^j\, H_{ij}^u
    \;+\; q_L^i\, d_R^j\, H_{ij}^d
  \Bigr)\,,
\end{equation}
where $i,j = 1,2,3$ are generation indices.  The effective Yukawa
matrices are~\cite{CacciapagliaSannino2024}:
\begin{equation}\label{eq:eff-yukawa}
  Y_{ij}^u \;=\; y\, \frac{\langle H_{ij}^u \rangle}{v}\,,
  \qquad
  Y_{ij}^d \;=\; y\, \frac{\langle H_{ij}^d \rangle}{v}\,,
\end{equation}
with the VEV sum rule:
\begin{equation}\label{eq:vev-sum}
  v^2 \;=\; \sum_{i,j}
  \Bigl(
    \langle H_{ij}^u \rangle^2 + \langle H_{ij}^d \rangle^2
  \Bigr)
  \;=\; \frac{1}{2}\,(246\;\text{GeV})^2\,.
\end{equation}
\textbf{Key insight:} The fermion mass hierarchy arises \emph{not} from
a hierarchy in the Yukawa couplings (they are all equal to $y$) but
from a hierarchy in the Higgs VEVs:
$\langle H_{33}^u \rangle \gg \langle H_{ij}^{u,d} \rangle$ for
$(i,j) \neq (3,3)$.  This is the ``scalar democracy''
mechanism~\cite{HillMachado2019}, which will be developed further in
Lecture~8.


%% ========================================================================
%%  SECTION 4: The Electric Theory
%% ========================================================================
\section{The Electric Theory}
\label{sec:electric-theory}

Having identified the magnetic spectrum, we now present the electric
theory that, \emph{if the duality holds}, provides the UV-complete
description of the SM at energies above the duality scale $\Lambda_D$.

\subsection{Electric field content for $n_g = 3$}

For three generations, the electric gauge group is
$\SU{2n_g - 4} = \SU{2}$ (from~\S\ref{ssec:generation-bound}).  The
electric theory contains~\cite{Sannino2011}:

\begin{center}
\renewcommand{\arraystretch}{1.3}
\begin{tabular}{c c c c c c}
\toprule
\multicolumn{6}{c}{\textbf{Electric theory: $\SU{2}$ with $\Nf = 6$}} \\
\midrule
Fields & $\SU{2}_{\text{el}}$ & $\SU{6}_L$ & $\SU{6}_R$ & $\UU{1}_V$ & $\UU{1}_{AF}$ \\
\midrule
$\lambda$ & Adj ($= \mathbf{3}$) & $\mathbf{1}$ & $\mathbf{1}$ & $0$ & $1$ \\
$p$       & $F$ ($= \mathbf{2}$) & $F$ & $\mathbf{1}$ & $1$  & $-1/2$ \\
$\widetilde{p}$ & $\bar{F}$ ($= \mathbf{2}$) & $\mathbf{1}$ & $\bar{F}$ & $-1$ & $-1/2$ \\
\midrule
$L$   & $\mathbf{1}$ & $F$ & $\mathbf{1}$ & $-3$ & $0$ \\
$\widetilde{L}$ & $\mathbf{1}$ & $\mathbf{1}$ & $\bar{F}$ & $3$ & $0$ \\
\bottomrule
\end{tabular}
\captionof{table}{Electric theory for $n_g = 3$: $\SU{2}_{\text{el}}$
with 6~fundamental Weyl fermions $p$, 6~antifundamental
$\widetilde{p}$, and one adjoint $\lambda$.
For $\SU{2}$, the fundamental and antifundamental are equivalent
($\bar{F} \cong F$).}
\label{tab:electric-ng3}
\end{center}

\begin{remark}[No elementary scalars in the electric theory]
  The electric theory of Table~\ref{tab:electric-ng3} contains
  \textbf{only fermions and gauge bosons}---no elementary scalars.
  \emph{If the duality holds}, this means that \textbf{the hierarchy
  problem does not exist} in the fundamental description: there are no
  scalar masses to protect, because there are no fundamental scalars.
  The Higgs doublets of the SM are, within the conjectured dual map, composites
  $\Phi_H = \langle \widetilde{Q}\lambda\lambda Q \rangle$
  (Eq.~\eqref{eq:composites}), and their masses are set by the strong
  dynamics of the electric theory at the duality scale $\Lambda_D$.
\end{remark}

\subsection{Self-duality of $\SU{2}$}

Note that for $n_g = 3$, the electric gauge group is $\SU{2}$ with
$\Nf = 6$ Dirac flavours, and the magnetic gauge group is
$\SU{X} = \SU{\Nf - N} = \SU{6 - 2} = \SU{4}$.  The magnetic $\SU{4}$
is the \emph{Pati--Salam} gauge group containing QCD $\SU{3}_c$ as a
subgroup.  The duality connects:
\[
  \underbrace{\SU{2}_{\text{el}}}_{\text{electric}}
  \quad \xleftrightarrow{\text{duality}} \quad
  \underbrace{\SU{4}_{\text{PS}} \supset \SU{3}_c}_{\text{magnetic
  (contains SM)}}\,.
\]
Note also that $\SU{2}$ is \emph{self-dual} under Seiberg duality with
$\Nf = 4$ flavours (recalled from Lecture~3).  Here we have $\Nf = 6$,
so the electric and magnetic theories have different gauge groups---the
duality is a \emph{non-trivial} map, not a self-duality.

\subsection{Energy hierarchy}

\emph{If the duality holds}, the physics is characterised by a hierarchy
of energy scales~\cite{CacciapagliaSannino2024}:
\begin{enumerate}
  \item $E < \Lambda_S \sim \text{few TeV}$: only SM degrees of freedom
    are active.
  \item $\Lambda_S < E < \Lambda_D$: the full magnetic spectrum appears
    (gaugino, squarks, 17 heavy Higgs doublets, mesino).
  \item $E > \Lambda_D \sim 10^{11}\;\text{GeV}$: the electric theory
    $\SU{2}_{\text{el}}$ becomes weakly coupled; this is the
    ``electric SM.''
  \item Optionally, $E > \Lambda_{\text{eGUT}}$: if the three gauge
    couplings of the electric theory unify, a grand unified theory
    emerges above $\Lambda_{\text{eGUT}}$.
\end{enumerate}
At the duality scale $\Lambda_D$, the electric and magnetic gauge
couplings are related by~\cite{CacciapagliaSannino2024}:
\begin{equation}\label{eq:coupling-match}
  \frac{1}{\alpha_e}
  \;\sim\; \alpha_m
  \;=\; \frac{g_m^2}{4\pi}\,.
\end{equation}
Strong coupling in one description corresponds to weak coupling in the
other---the defining feature of duality (recall Lecture~3,
\S\ref{sec:general-duality}).


%% ========================================================================
%%  SECTION 5: Anomaly Matching for the Full SM Gauge Group
%% ========================================================================
\section{Anomaly Matching for $\SU{3}_c \times \SU{2}_L \times \UU{1}_Y$}
\label{sec:anomaly-matching}

The central non-trivial consistency check of the dualised SM
construction is 't~Hooft anomaly matching (Lecture~1) between the
electric and magnetic descriptions.  We now compute all independent
anomaly conditions for the gauged subgroup
$\SU{3}_c \times \SU{2}_L \times \UU{1}_Y$ of the global symmetry.

\begin{remark}[Anomaly matching as consistency check]
  Recall from Lecture~1: 't~Hooft anomaly matching constrains the
  massless fermion spectrum by requiring that all anomaly coefficients
  for global symmetries match between UV and IR descriptions.  Here, we
  treat $\SU{3}_c \times \SU{2}_L \times \UU{1}_Y$ as the gauged
  symmetry whose anomalies must \emph{cancel} (gauge anomaly
  cancellation), and we verify that both the electric and magnetic
  theories satisfy the same anomaly conditions.
\end{remark}

\subsection{Independent anomaly conditions}

For a gauge group $\SU{3}_c \times \SU{2}_L \times \UU{1}_Y$, the
independent anomaly conditions for gauge anomaly cancellation
are~\cite{tHooft1980}:

\begin{enumerate}[label=(\roman*)]
  \item $\SU{3}^3$: the cubic $\SU{3}_c$ anomaly.
  \item $\SU{2}^3$: the cubic $\SU{2}_L$ anomaly.
  \item $\SU{3}^2\,\UU{1}_Y$: the mixed colour--hypercharge anomaly.
  \item $\SU{2}^2\,\UU{1}_Y$: the mixed weak--hypercharge anomaly.
  \item $\UU{1}_Y^3$: the cubic hypercharge anomaly.
  \item $\UU{1}_Y$--gravitational: the mixed hypercharge--gravitational
    anomaly.
\end{enumerate}

\subsection{Condition (ii): $\SU{2}^3$ vanishes trivially}

The cubic Casimir $d_{abc}$ vanishes for $\SU{2}$:
\begin{equation}\label{eq:su2-cubic}
  A_{\SU{2}^3} \;=\; \sum_{\text{Weyl}} \Tr\bigl[T^a_{\SU{2}} \{T^b_{\SU{2}}, T^c_{\SU{2}}\}\bigr]
  \;=\; 0
  \qquad\text{(all representations)}\,.
\end{equation}
This is a group-theoretic identity for $\SU{2}$ and provides no
constraint on the spectrum.

\subsection{Computing the anomaly coefficients}

For the remaining five non-trivial conditions, we need the SM quantum
numbers of all magnetic fermions.  Using the flavour
decomposition~\eqref{eq:flavour-decomp} with $n_g = 3$ and the
hypercharge formula~\eqref{eq:hypercharge}, the magnetic
\textbf{fermion} content (only fermions contribute to anomalies)
decomposes as follows.

\medskip

\textbf{Magnetic fermions and their SM quantum numbers:}

\renewcommand{\arraystretch}{1.4}
\begin{center}
\begin{tabular}{l c c c c c}
\toprule
\textbf{Field} & $\SU{3}_c$ & $\SU{2}_L$ & $Y$
  & \textbf{Weyl count} & \textbf{Type} \\
\midrule
\multicolumn{6}{c}{\emph{SM quarks (from $q$, $\widetilde{q}$):}} \\
$q_L^i$ & $\mathbf{3}$ & $\mathbf{2}$ & $+1/6$ & $3 \times 2 = 6$ & Weyl \\
$u_R^i$ & $\mathbf{3}$ & $\mathbf{1}$ & $+2/3$ & $3$ & Weyl \\
$d_R^i$ & $\mathbf{3}$ & $\mathbf{1}$ & $-1/3$ & $3$ & Weyl \\
$\widetilde{q}_L^i$ & $\bar{\mathbf{3}}$ & $\mathbf{2}$ & $-1/6$ & $6$ & Weyl \\
$\widetilde{u}_R^i$ & $\bar{\mathbf{3}}$ & $\mathbf{1}$ & $-2/3$ & $3$ & Weyl \\
$\widetilde{d}_R^i$ & $\bar{\mathbf{3}}$ & $\mathbf{1}$ & $+1/3$ & $3$ & Weyl \\
\midrule
\multicolumn{6}{c}{\emph{SM leptons (from $l$, $\widetilde{l}$):}} \\
$\ell_L^i$ & $\mathbf{1}$ & $\mathbf{2}$ & $-1/2$ & $6$ & Weyl \\
$\nu_R^i$ & $\mathbf{1}$ & $\mathbf{1}$ & $0$ & $3$ & Weyl \\
$e_R^i$ & $\mathbf{1}$ & $\mathbf{1}$ & $-1$ & $3$ & Weyl \\
$\widetilde{\ell}_L^i$ & $\mathbf{1}$ & $\mathbf{2}$ & $+1/2$ & $6$ & Weyl \\
$\widetilde{\nu}_R^i$ & $\mathbf{1}$ & $\mathbf{1}$ & $0$ & $3$ & Weyl \\
$\widetilde{e}_R^i$ & $\mathbf{1}$ & $\mathbf{1}$ & $+1$ & $3$ & Weyl \\
\midrule
\multicolumn{6}{c}{\emph{Gaugino (from $\lambda_m$):}} \\
$\lambda_m$ & $\mathbf{8}$ & $\mathbf{1}$ & $0$ & $8$ & Weyl \\
\midrule
\multicolumn{6}{c}{\emph{Mesino (from $M$):}} \\
$M_{ij}$ & $\mathbf{1}$ & various & various & $36$ & Weyl \\
\bottomrule
\end{tabular}
\captionof{table}{SM quantum numbers of all magnetic Weyl fermions.
Scalars ($\phi$, $\widetilde{\phi}$, $\Phi_H$) do \emph{not}
contribute to anomalies.}
\label{tab:magnetic-SM-quantum}
\end{center}

\medskip

A crucial observation is that the magnetic fermion content can be
decomposed into two sectors:
\begin{enumerate}
  \item \textbf{Vector-like pairs:} The SM quarks and leptons come in
    left-right pairs ($q$, $\widetilde{q}$) and ($l$, $\widetilde{l}$)
    with opposite $\UU{1}_V$ charges and conjugate gauge
    representations.  Each left-right pair contributes zero net
    anomaly to any purely gauge anomaly condition.
  \item \textbf{Extra states:} The gaugino $\lambda_m$ and mesino $M$
    have specific $\SU{3}_c$, $\SU{2}_L$, and $\UU{1}_Y$ quantum
    numbers.
\end{enumerate}

\subsection{Condition (i): $\SU{3}^3$}

The cubic $\SU{3}_c$ anomaly coefficient is:
\begin{equation}\label{eq:su3-cubed}
  A_{\SU{3}^3}
  \;=\; \sum_{\text{Weyl}} A(R_{\text{colour}})\,,
\end{equation}
where $A(\mathbf{3}) = 1$, $A(\bar{\mathbf{3}}) = -1$,
$A(\mathbf{8}) = 0$ (adjoint is real), and $A(\mathbf{1}) = 0$.

\textbf{SM quarks:} Per generation, the SM quark contribution is
$A(q_L) + A(u_R) + A(d_R) = 2 \cdot 1 + 1 + 1 = 4$ from
left-handed quarks (two $\SU{2}_L$ doublet components) and
$A(\widetilde{q}_L) + A(\widetilde{u}_R) + A(\widetilde{d}_R)
= 2 \cdot (-1) + (-1) + (-1) = -4$ from the conjugate quarks.
Total per generation: $4 + (-4) = 0$.

\textbf{Extra states:} $A(\lambda_m) = A(\mathbf{8}) = 0$ (adjoint is
real).  The mesino $M$ is a colour singlet: $A(\mathbf{1}) = 0$.

\begin{equation}\label{eq:su3-result}
  \boxed{A_{\SU{3}^3} \;=\; 0\,.}
  \qquad \text{(Vector-like pairs cancel generation by generation.)}
\end{equation}

\subsection{Condition (iii): $\SU{3}^2\,\UU{1}_Y$}

This anomaly coefficient sums $Y \cdot T(R_{\text{colour}})$ over all
colour-charged Weyl fermions:
\begin{equation}\label{eq:su3sq-u1}
  A_{\SU{3}^2\,\UU{1}_Y}
  \;=\; \sum_{\text{colour-charged Weyl}} Y \cdot T(R_{\text{colour}})\,.
\end{equation}
With $T(\mathbf{3}) = T(\bar{\mathbf{3}}) = 1/2$ and
$T(\mathbf{8}) = 3$:

\textbf{SM quarks (per generation):}
\begin{align}
  &q_L:\; 2 \times \tfrac{1}{6} \times \tfrac{1}{2} = \tfrac{1}{6}\,,
  \quad
  u_R:\; \tfrac{2}{3} \times \tfrac{1}{2} = \tfrac{1}{3}\,,
  \quad
  d_R:\; (-\tfrac{1}{3}) \times \tfrac{1}{2} = -\tfrac{1}{6}\,,
  \nonumber\\
  &\text{Sum} = \tfrac{1}{6} + \tfrac{1}{3} - \tfrac{1}{6} = \tfrac{1}{3}\,.
  \label{eq:su3sq-quarks}
\end{align}
\textbf{Conjugate quarks (per generation):}
\begin{align}
  &\widetilde{q}_L:\; 2 \times (-\tfrac{1}{6}) \times \tfrac{1}{2} = -\tfrac{1}{6}\,,
  \quad
  \widetilde{u}_R:\; (-\tfrac{2}{3}) \times \tfrac{1}{2} = -\tfrac{1}{3}\,,
  \quad
  \widetilde{d}_R:\; \tfrac{1}{3} \times \tfrac{1}{2} = \tfrac{1}{6}\,,
  \nonumber\\
  &\text{Sum} = -\tfrac{1}{6} - \tfrac{1}{3} + \tfrac{1}{6} = -\tfrac{1}{3}\,.
  \label{eq:su3sq-conj-quarks}
\end{align}
\textbf{Quarks total (3 generations):} $3 \times (\tfrac{1}{3} - \tfrac{1}{3}) = 0$.

\textbf{Gaugino:} $\lambda_m$ has $Y = 0$ (colour octet, $\SU{2}_L$
singlet, $Q_V = 0$), so $Y \cdot T(\mathbf{8}) = 0 \times 3 = 0$.

\begin{equation}\label{eq:su3sq-u1-result}
  \boxed{A_{\SU{3}^2\,\UU{1}_Y} \;=\; 0\,.}
  \qquad\text{(Vector-like quark pairs cancel; gaugino is neutral.)}
\end{equation}

\subsection{Condition (iv): $\SU{2}^2\,\UU{1}_Y$}

This anomaly coefficient sums $Y \cdot T(R_{\text{weak}})$ over all
$\SU{2}_L$ doublet Weyl fermions:
\begin{equation}\label{eq:su2sq-u1}
  A_{\SU{2}^2\,\UU{1}_Y}
  \;=\; \sum_{\SU{2}_L\text{-doublet Weyl}} Y \cdot T(\mathbf{2})\,.
\end{equation}
With $T(\mathbf{2}) = 1/2$ for the $\SU{2}_L$ fundamental:

\textbf{SM quark doublets (per generation):}
$q_L:\; Y = +1/6$, $T = 1/2$, colour multiplicity $3$:
contribution $= 3 \times \tfrac{1}{6} \times \tfrac{1}{2} = \tfrac{1}{4}$.

\textbf{Conjugate quark doublets (per generation):}
$\widetilde{q}_L:\; Y = -1/6$, $T = 1/2$, colour multiplicity $3$:
contribution $= 3 \times (-\tfrac{1}{6}) \times \tfrac{1}{2} = -\tfrac{1}{4}$.

\textbf{Lepton doublets (per generation):}
$\ell_L:\; Y = -1/2$, $T = 1/2$, colour multiplicity $1$:
contribution $= 1 \times (-\tfrac{1}{2}) \times \tfrac{1}{2} = -\tfrac{1}{4}$.

\textbf{Conjugate lepton doublets (per generation):}
$\widetilde{\ell}_L:\; Y = +1/2$, $T = 1/2$, colour multiplicity $1$:
contribution $= 1 \times \tfrac{1}{2} \times \tfrac{1}{2} = \tfrac{1}{4}$.

\textbf{Total per generation:}
$\tfrac{1}{4} - \tfrac{1}{4} - \tfrac{1}{4} + \tfrac{1}{4} = 0$.

\textbf{Gaugino:} $\SU{2}_L$ singlet---does not contribute.

\textbf{Mesino:} The mesino $M$ transforms under $\SU{6}_L \times \SU{6}_R$.
Under the flavour decomposition~\eqref{eq:flavour-decomp}, $M$ contains
components that are $\SU{2}_L$ doublets with various hypercharges.
However, $M$ transforms as $(F, \bar{F})$ under
$\SU{6}_L \times \SU{6}_R$ with $Q_V = 0$.
Its $\SU{2}_L$ doublet components come in conjugate pairs (from the
$\SU{2}_L \subset \SU{6}_L$ and $\SU{2}_R \subset \SU{6}_R$
decomposition), which contribute with opposite signs to the anomaly.
The mesino's net contribution to $\SU{2}^2\,\UU{1}_Y$ vanishes.

\begin{equation}\label{eq:su2sq-u1-result}
  \boxed{A_{\SU{2}^2\,\UU{1}_Y} \;=\; 0\,.}
  \qquad\text{(All contributions cancel pairwise.)}
\end{equation}

\subsection{Condition (v): $\UU{1}_Y^3$}

The cubic hypercharge anomaly sums $Y^3$ over all Weyl fermions:
\begin{equation}\label{eq:u1-cubed}
  A_{\UU{1}_Y^3}
  \;=\; \sum_{\text{all Weyl}} \dim(R_{\text{colour}}) \cdot \dim(R_{\text{weak}}) \cdot Y^3\,.
\end{equation}

\textbf{SM quarks (per generation):}
\begin{align}
  &q_L:\; 3 \times 2 \times (1/6)^3 = 6/216 = 1/36\,,
  \nonumber\\
  &u_R:\; 3 \times 1 \times (2/3)^3 = 3 \times 8/27 = 24/27 = 8/9\,,
  \nonumber\\
  &d_R:\; 3 \times 1 \times (-1/3)^3 = 3 \times (-1/27) = -3/27 = -1/9\,.
  \label{eq:u1-quarks}
\end{align}
\textbf{Conjugate quarks (per generation):}
\begin{align}
  &\widetilde{q}_L:\; 3 \times 2 \times (-1/6)^3 = -1/36\,,
  \nonumber\\
  &\widetilde{u}_R:\; 3 \times 1 \times (-2/3)^3 = -8/9\,,
  \nonumber\\
  &\widetilde{d}_R:\; 3 \times 1 \times (1/3)^3 = 1/9\,.
  \label{eq:u1-conj-quarks}
\end{align}
Quark + conjugate quark total per generation: 0 (each term cancels its
partner).

\textbf{SM leptons (per generation):}
\begin{align}
  &\ell_L:\; 1 \times 2 \times (-1/2)^3 = -1/4\,,
  \nonumber\\
  &\nu_R:\; 1 \times 1 \times 0^3 = 0\,,
  \nonumber\\
  &e_R:\; 1 \times 1 \times (-1)^3 = -1\,.
  \label{eq:u1-leptons}
\end{align}
\textbf{Conjugate leptons (per generation):}
\begin{align}
  &\widetilde{\ell}_L:\; 1 \times 2 \times (1/2)^3 = 1/4\,,
  \nonumber\\
  &\widetilde{\nu}_R:\; 1 \times 1 \times 0^3 = 0\,,
  \nonumber\\
  &\widetilde{e}_R:\; 1 \times 1 \times (1)^3 = 1\,.
  \label{eq:u1-conj-leptons}
\end{align}
Lepton + conjugate lepton total per generation: 0.

\textbf{Gaugino:} $Y = 0$, so $Y^3 = 0$.

\textbf{Mesino:} $M$ has $Q_V = 0$, and its $\SU{2}_R$ components come
in pairs with $T_R^3 = \pm 1/2$ and $T_R^3 = 0$.  Since $Q_V = 0$,
$Y = T_R^3$ for the mesino.  The $T_R^3 = +1/2$ and $T_R^3 = -1/2$
components are related by the $\SU{2}_R$ structure and have equal
multiplicities, so their $Y^3$ contributions cancel:
$(+1/2)^3 + (-1/2)^3 = 0$.  The $T_R^3 = 0$ components contribute
$0^3 = 0$.

\begin{equation}\label{eq:u1-cubed-result}
  \boxed{A_{\UU{1}_Y^3} \;=\; 0\,.}
  \qquad\text{(Vector-like pairs and $\SU{2}_R$ pairing enforce cancellation.)}
\end{equation}

\subsection{Condition (vi): $\UU{1}_Y$--gravitational}

The gravitational anomaly sums $Y$ over all Weyl fermions:
\begin{equation}\label{eq:u1-grav}
  A_{\UU{1}_Y\text{-grav}}
  \;=\; \sum_{\text{all Weyl}} \dim(R_{\text{colour}}) \cdot \dim(R_{\text{weak}}) \cdot Y\,.
\end{equation}

\textbf{SM quarks (per generation):}
\begin{align}
  &q_L:\; 3 \times 2 \times 1/6 = 1\,,
  \nonumber\\
  &u_R:\; 3 \times 1 \times 2/3 = 2\,,
  \nonumber\\
  &d_R:\; 3 \times 1 \times (-1/3) = -1\,.
  \nonumber
\end{align}
\textbf{Conjugate quarks (per generation):}
\begin{align}
  &\widetilde{q}_L:\; 3 \times 2 \times (-1/6) = -1\,,
  \nonumber\\
  &\widetilde{u}_R:\; 3 \times 1 \times (-2/3) = -2\,,
  \nonumber\\
  &\widetilde{d}_R:\; 3 \times 1 \times (1/3) = 1\,.
  \nonumber
\end{align}
Quark + conjugate total per generation: $(1+2-1) + (-1-2+1) = 0$.

\textbf{SM leptons (per generation):}
$\ell_L: 2 \times (-1/2) = -1$; $\nu_R: 0$; $e_R: -1$.
\textbf{Conjugate:}
$\widetilde{\ell}_L: 2 \times (+1/2) = +1$;
$\widetilde{\nu}_R: 0$; $\widetilde{e}_R: +1$.
Total per generation: $(-1+0-1) + (1+0+1) = 0$.

\textbf{Gaugino:} $Y = 0$, so contribution $= 8 \times 1 \times 0 = 0$.

\textbf{Mesino:} By the same $\SU{2}_R$ pairing argument as in
Condition~(v), the mesino contributions cancel:
$(+1/2) + (-1/2) = 0$ for each doublet pair, and $0$ for singlets.

\begin{equation}\label{eq:u1-grav-result}
  \boxed{A_{\UU{1}_Y\text{-grav}} \;=\; 0\,.}
  \qquad\text{(Vector-like structure ensures vanishing.)}
\end{equation}

\subsection{Summary of anomaly matching}

\begin{center}
\renewcommand{\arraystretch}{1.4}
\begin{tabular}{l c c c}
\toprule
\textbf{Anomaly condition}
  & \textbf{Magnetic (SM + partners)}
  & \textbf{Electric}
  & \textbf{Status} \\
\midrule
$\SU{3}^3$ & $0$ & $0$ & \checkmark \\
$\SU{2}^3$ & $0$ (trivial: $d_{abc} = 0$) & $0$ & \checkmark \\
$\SU{3}^2\,\UU{1}_Y$ & $0$ & $0$ & \checkmark \\
$\SU{2}^2\,\UU{1}_Y$ & $0$ & $0$ & \checkmark \\
$\UU{1}_Y^3$ & $0$ & $0$ & \checkmark \\
$\UU{1}_Y\text{-grav}$ & $0$ & $0$ & \checkmark \\
\bottomrule
\end{tabular}
\captionof{table}{Summary of anomaly matching for
$\SU{3}_c \times \SU{2}_L \times \UU{1}_Y$.  All six conditions
are satisfied.}
\label{tab:anomaly-summary}
\end{center}

The key mechanism is the \textbf{vector-like structure} of the magnetic
theory: quarks and conjugate quarks ($q$, $\widetilde{q}$) carry
opposite quantum numbers, as do leptons and conjugate leptons ($l$,
$\widetilde{l}$).  The extra states (gaugino, mesino) either are
neutral under $\UU{1}_Y$ (gaugino) or come in $\SU{2}_R$ pairs that
cancel (mesino).

\subsection{Electric side: why the anomalies also vanish}

In the electric theory, $\SU{3}_c \times \SU{2}_L \times \UU{1}_Y$
is a \emph{subgroup of the global symmetry} (it is gauged only at
lower energies, in the magnetic description).  The 't~Hooft anomaly
coefficients of this subgroup must match between the electric and
magnetic descriptions.

The electric fermions (Table~\ref{tab:electric-ng3}) decompose under
$\SU{3}_c \times \SU{2}_L \times \UU{1}_Y$ using the same
embedding~\eqref{eq:flavour-decomp} and
hypercharge~\eqref{eq:hypercharge}:
\begin{itemize}
  \item $Q$ ($\SU{3}$ triplet, $\SU{6}_L$ fundamental, $Q_V = +1$):
    decomposes into 3~generations of $\SU{2}_L$ doublets with
    $Y = +1/6$, each carrying an $\SU{2}_{\text{el}}$ fundamental
    index (multiplicity~2).
  \item $\widetilde{Q}$ ($\bar{\mathbf{3}}$, $\SU{6}_R$
    antifundamental, $Q_V = -1$): decomposes into 3~generations of
    $\SU{2}_R$ doublets.  The $T_R^3 = +1/2$ component has
    $Y = 1/2 - 1/6 = 1/3$ and the $T_R^3 = -1/2$ component has
    $Y = -1/2 - 1/6 = -2/3$.  Each carries
    $\SU{2}_{\text{el}}$ antifundamental (multiplicity~2).
  \item $\lambda$ ($\SU{3}$ octet, flavour singlet, $Q_V = 0$,
    $Y = 0$): carries $\SU{2}_{\text{el}}$ adjoint (multiplicity~3).
  \item $L$ ($\SU{3}$ singlet, $\SU{6}_L$ fundamental, $Q_V = -3$):
    decomposes into 3~generations of $\SU{2}_L$ doublets with
    $Y = -1/2$.  These are $\SU{2}_{\text{el}}$ singlets.
  \item $\widetilde{L}$ ($\SU{3}$ singlet, $\SU{6}_R$
    antifundamental, $Q_V = +3$): $\SU{2}_R$ doublets with
    $Y = +1$ and $Y = 0$.  $\SU{2}_{\text{el}}$ singlets.
\end{itemize}
Crucially, the $\SU{3}_c \times \SU{2}_L \times \UU{1}_Y$ quantum
numbers of the electric fermions are \textbf{identical} to those of
the corresponding magnetic fermions (compare with
Table~\ref{tab:magnetic-SM-quantum})---the only difference is the
$\SU{2}_{\text{el}}$ multiplicity factor.  Since every anomaly
coefficient of a product group factorises as
$A = \dim(R_{\text{el}}) \times A_{\text{SM}}$, and $A_{\text{SM}} = 0$
for each condition (as computed above), the electric anomalies also
vanish:
\begin{equation}\label{eq:electric-anomalies}
  A_{\text{electric}} \;=\; \dim(R_{\SU{2}_{\text{el}}}) \times
  \underbrace{A_{\SU{3}_c \times \SU{2}_L \times \UU{1}_Y}}_{ = \,0}
  \;=\; 0\,.
\end{equation}
This confirms the anomaly matching between electric and magnetic
descriptions: $A_{\text{electric}} = A_{\text{magnetic}} = 0$ for all
six conditions.

\begin{warningbox}[title={Anomaly matching is necessary but not sufficient}]
  The anomaly matching demonstrated in Table~\ref{tab:anomaly-summary}
  is a \textbf{necessary condition} for the conjectured duality between
  the electric $\SU{2}$ theory and the magnetic SM, not a
  \textbf{sufficient} one.  It constrains the field content of the
  magnetic theory but does not prove the dynamical equivalence of the
  two descriptions.

  \medskip

  \emph{If the anomaly matching had failed}---if any coefficient were
  non-zero on one side and zero on the other---the duality as stated
  would be \textbf{inconsistent} and the field content would need
  modification.  The fact that all conditions are satisfied is
  \textbf{evidence} for the construction's consistency, but the
  construction remains \textbf{conjectural} (recall the central
  vulnerability of Lecture~6, \S\ref{sec:vulnerability}).
\end{warningbox}

\subsection{Recovery of SM anomaly cancellation}

As a cross-check, we verify that when the extra magnetic states are
``removed'' (their contributions set to zero), the standard SM anomaly
cancellation is recovered.

The extra magnetic fermions are:
\begin{itemize}
  \item Gaugino $\lambda_m$: $Y = 0$ for all anomaly conditions.
    Removing it changes nothing.
  \item Mesino $M$: contributes zero to all anomaly conditions (by the
    $\SU{2}_R$ pairing argument).  Removing it changes nothing.
  \item Conjugate quarks $\widetilde{q}$ and conjugate leptons
    $\widetilde{l}$: these form exact vector-like pairs with $q$ and
    $l$.  Removing both $\widetilde{q}$ and $\widetilde{l}$ leaves
    exactly the SM quark and lepton content.
\end{itemize}

After removing all extra states, the remaining fields are exactly three
generations of SM quarks and leptons:
$\{q_L, u_R, d_R, \ell_L, \nu_R, e_R\} \times 3$, and the standard
SM anomaly cancellation from Lecture~1 is recovered.  This consistency
check confirms that the anomaly matching of the full magnetic theory
is a \emph{refinement} of the standard SM result, not a replacement.


%% ========================================================================
%%  SUMMARY
%% ========================================================================
\section{Summary}
\label{sec:summary-lec7}

\begin{tcolorbox}[
  enhanced,
  colback=green!5!white,
  colframe=green!60!black,
  fonttitle=\bfseries,
  title={Lecture 7: Summary}
]
\begin{enumerate}
  \item \textbf{General duality framework:}
    An $\SU{N}$ gauge theory with $\Nf$ Dirac fundamentals and one
    adjoint Weyl fermion has, \emph{if the conjectured non-SUSY duality
    holds}, a magnetic dual $\SU{X}$ with $X = \Nf - N$ and a
    quasi-SUSY spectrum: gaugino, magnetic quarks, squarks, mesino,
    and mes-Higgs~(\S\ref{sec:general-duality}).

  \item \textbf{Pati--Salam embedding:}
    Embedding quarks and leptons in $\SU{4}$ (lepton~$=$~fourth colour,
    Pati--Salam~\cite{PS1974}), the flavour symmetry decomposes as
    $\SU{\Nf}_{L/R} \supset \SU{n_g} \times \SU{2}_{L/R}$, and
    hypercharge is identified as $Y = T_R^3 + Q_V/6$
    (\S\ref{sec:pati-salam}).  All six SM hypercharges are reproduced
    exactly.

  \item \textbf{Generation bound:}
    $2n_g - 4 \geq 2$ gives $n_g \geq 3$.  \emph{If the duality
    holds}, three or more generations are a mathematical necessity
    (\S\ref{ssec:generation-bound}).

  \item \textbf{Magnetic spectrum:}
    For $n_g = 3$, the magnetic theory contains all SM fields plus
    quasi-SUSY partners (gaugino, squarks, mesino) and 18 Higgs
    doublets (9~bi-doublets $\times$ 2), of which one acquires the
    electroweak VEV~(\S\ref{sec:magnetic-spectrum}).

  \item \textbf{Electric theory:}
    The UV description is $\SU{2}$ with 6 Dirac flavours and an
    adjoint---\emph{no elementary scalars}.  If the duality holds, all scalars in the
    magnetic theory are composites
    (\S\ref{sec:electric-theory}).

  \item \textbf{Anomaly matching:}
    All six independent anomaly conditions for
    $\SU{3}_c \times \SU{2}_L \times \UU{1}_Y$ are satisfied.
    The vector-like structure of the magnetic theory ensures
    cancellation.  This is necessary but not sufficient for the
    duality~(\S\ref{sec:anomaly-matching}).
\end{enumerate}
\end{tcolorbox}

\begin{previewbox}[title={Preview: Lecture 8}]
  In Lecture~8, we will develop the quantitative consequences of the
  dualised SM:
  \begin{itemize}
    \item The scalar democracy mechanism: universal Yukawa coupling $y$
      combined with hierarchical Higgs VEVs produces the observed
      fermion mass hierarchy.
    \item Scale matching: the duality scale
      $\Lambda_D \sim 10^{11}\;\text{GeV}$, the Fermi scale
      $v \approx 246\;\text{GeV}$, and
      $\Lambda_{\QCD} \approx 200\;\text{MeV}$ from the dual
      description.
    \item Collider signatures: the quasi-SUSY spectrum at the TeV scale.
  \end{itemize}

  \textbf{Every prediction derived in Lecture~8 is conditional on
  the conjectured non-SUSY duality being correct.}
\end{previewbox}


%% ========================================================================
%%  BIBLIOGRAPHY
%% ========================================================================
\begin{thebibliography}{99}

\bibitem{CacciapagliaSannino2024}
  G.~Cacciapaglia and F.~Sannino,
  ``Charting Standard Model duality and its signatures,''
  \textit{Phys.~Rev.} \textbf{D111} (2025) 035013
  [\texttt{arXiv:2407.17281}].

\bibitem{Sannino2011}
  F.~Sannino,
  ``The Standard Model is natural as magnetic gauge theory,''
  \textit{Mod.~Phys.~Lett.} \textbf{A26} (2011) 1763
  [\texttt{arXiv:1102.5100}].

\bibitem{PS1974}
  J.~C.~Pati and A.~Salam,
  ``Lepton number as the fourth `color',''
  \textit{Phys.~Rev.} \textbf{D10} (1974) 275.

\bibitem{HillMachado2019}
  C.~T.~Hill, P.~A.~N.~Machado, A.~E.~Thomsen, and J.~Turner,
  ``Scalar democracy,''
  \textit{Phys.~Rev.} \textbf{D100} (2019) 015015
  [\texttt{arXiv:1902.07214}].

\bibitem{AMPS2011}
  O.~Antipin, M.~Mojaza, C.~Pica, and F.~Sannino,
  ``Magnetic fixed points and emergent supersymmetry,''
  \textit{JHEP} \textbf{1306} (2013) 037
  [\texttt{arXiv:1105.1510}].

\bibitem{tHooft1980}
  G.~'t~Hooft,
  ``Naturalness, chiral symmetry, and spontaneous chiral symmetry
  breaking,'' in \textit{Recent Developments in Gauge Theories},
  NATO ASI Series~B59 (Plenum, 1980) 135.

\bibitem{Seiberg1995}
  N.~Seiberg,
  ``Electric--magnetic duality in supersymmetric non-Abelian gauge
  theories,''
  \textit{Nucl.~Phys.} \textbf{B435} (1995) 129
  [\texttt{hep-th/9411149}].

\bibitem{CST2011}
  C.~Csaki, Y.~Shirman, and J.~Terning,
  ``A Seiberg dual for the MSSM: partially composite $W$ and $Z$,''
  \textit{Phys.~Rev.} \textbf{D84} (2011) 095011
  [\texttt{arXiv:1106.3074}].

\end{thebibliography}

\end{document}
